\chapter{Фазово-взвешенный третий момент и амплитудный запрет}
\label{ch:v6-single-thread-phase-weighted-moment-lift-gate}

\section{Постановка}

Предыдущая глава закрыла общую сырую моментную плотность, но оставила
минимальную возможность: поле $Q$ измеряет обычную плотность локальных
проходов, а поле $T$ получает на тех же проходах канонические фазы
остаточной $Z_3$-голономии или Real-пары.

Такое различение согласуется с общей литературной границей: связность
длинной цепи накладывает дополнительные законы сохранения на её
ориентационные моменты
\cite{single-thread-moment-svensek2013,single-thread-moment-svensek2016},
а тетраэдрический третий момент является самостоятельным параметром
порядка
\cite{single-thread-moment-radzihovsky2000,single-thread-moment-lubensky2002}.
Однако фазовый подъём допустим только в том случае, если уже существующие
характеры воспроизводят $T_*$ без нового модуля.

\section{Линейно независимые тетраэдрические вклады}

Обозначим бесследовые симметрические вклады четырёх направлений через
\begin{equation}
 C_a=\operatorname{STF}(n_a^{\otimes3}).
\end{equation}
Как элементы семимерного пространства спина три, четыре столбца $C_a$
имеют ранг четыре. Их сингулярные числа равны
\begin{equation}
 0.942809042,
 \qquad
 0.486864496^{(3)}.
\end{equation}
Следовательно, коэффициенты разложения по $C_a$ определяются однозначно.
Канонический тетраэдрический тензор имеет вид
\begin{equation}
 T_*=\sum_{a=1}^{4}C_a.
\end{equation}

Пусть упорядоченный второй момент по-прежнему фиксирует веса
\begin{equation}
 w=(p,q,q,q),
 \qquad
 p=0.9011874254,
 \qquad
 q=0.03293752488.
\end{equation}
Фазово-взвешенный кандидат равен
\begin{equation}
 T_\chi=\sum_aw_a\chi_aC_a,
 \qquad |\chi_a|=1.
\end{equation}

\begin{theorem}[Амплитудный запрет унитарного характера]
Равенство $T_\chi=cT_*$ возможно только при
$w_a\chi_a=c$ для всех четырёх направлений. Поэтому унитарный характер
$|\chi_a|=1$ требует равенства всех $w_a$. При $p\ne q$ никакой набор
фаз, включая все характеры $Z_3$, не воспроизводит $T_*$.
\end{theorem}

Это точный линейно-алгебраический запрет, а не результат ограниченного
численного поиска.

\section{Полный скан $Z_3$ и Real-пары}

Выбранная ось $n_1$ неподвижна, а три остальные оси образуют орбиту
остаточной $Z_3$. Проверены все перестановки характеров
$1,\omega,\omega^2$ на этой орбите и полное непрерывное вещественное
пространство, натянутое на действительную и мнимую части каждой
Real-комбинации.

Тривиальный характер воспроизводит прежний остаток
\begin{equation}
 \varepsilon_{\rm triv}=0.6133541143.
\end{equation}
Лучший нетривиальный характер имеет
\begin{equation}
 \varepsilon_{Z_3}=0.6671606493,
 \qquad
 |\cos(T_\chi,T_*)|=0.7449138662.
\end{equation}
Следовательно, фазовое скручивание не только не устраняет амплитудное
несовпадение, но в лучшем нетривиальном секторе увеличивает его.

Контрнаправленная Real-пара может обнулить физический первый момент и
менять действительные амплитуды посредством линейной комбинации
сопряжённых ветвей. Но неподвижная ось и одно направление трёхэлементной
орбиты несут один и тот же тривиальный фазовый множитель при разных весах
$p$ и $q$. Поэтому даже полный непрерывный Real-спан не снимает запрет
третьего момента.

\section{Требуемый, но не выведенный вес}

Точное выравнивание потребовало бы амплитуд
\begin{equation}
 h_a\propto\frac1{w_a},
\end{equation}
то есть
\begin{equation}
 h=(1.109647085,
 30.36050837,
 30.36050837,
 30.36050837)
\end{equation}
с отношением
\begin{equation}
 \frac{h_{\rm off}}{h_{\rm selected}}
 =\frac pq=27.36050837.
\end{equation}
Это не фаза и не дискретный характер. Формальный обратный вес напоминает
модулярный $R^{-1}$, но предыдущий гейт уже установил: обратное состояние
канонически входит в оператор модулярного потока, а не автоматически в
физическую энергию или наблюдаемое поле. Поэтому использовать $1/w_a$
вручную запрещено.

\begin{nogocriterion}
Нельзя объявлять отношение $27.36050837$ топологической кратностью или
связностным числом. Оно зависит от упорядоченного спектра и должно быть
выведено одним родительским оператором до сравнения с $T_*$.
\end{nogocriterion}

\section{Следующий различающий тест}

Чисто фазовый маршрут закрыт. Следующая минимальная возможность ---
связностно-взвешенный момент: локальная сшивка проходов может считать не
время пребывания в направлении, а частоту переходов между направлениями.

Нулевая гипотеза: уже существующие частичная изометрия, тетраэдрический
инцидентный оператор или смешанная кривизна $F_{QT}$ канонически создают
амплитудный вес, эквивалентный требуемому $1/w_a$, без нового параметра.

Стоп-критерий: если все существующие операторы дают только единичные
характеры, полиномиальные веса $w_a$ или свободную нормировку, прямой мост
от одной нити к совместному $Q+T$-вакууму закрывается.

\begin{protocol}
\begin{enumerate}
 \item ранг четырёх вкладов спина три: \textbf{четыре};
 \item точное фазовое решение при $p\ne q$: \textbf{невозможно};
 \item остаток тривиального характера: \textbf{$0.6133541143$};
 \item лучший остаток нетривиального $Z_3$: \textbf{$0.6671606493$};
 \item Real-пара снимает первый момент: \textbf{может};
 \item Real-пара снимает амплитудный запрет: \textbf{нет};
 \item требуемое отношение модулей: \textbf{$27.36050837$};
 \item чисто фазово-взвешенный подъём: \textbf{закрыт};
 \item одна глобальная нить: \textbf{не опровергнута};
 \item следующий тест: \textbf{связностно-взвешенный момент};
 \item рождение материи: \textbf{не закрыто}.
\end{enumerate}
\end{protocol}

Машинный сертификат сохранён в
\path{s2t_v6_single_thread_phase_weighted_moment_lift_gate_results.json}.